\chapter{Conhecendo o Paraguai}\label{conhecendo-o-paraguai}

\section{1. Paraguai em perspectiva}\label{paraguai-em-perspectiva}

\subsection{1.1 Por que o Paraguai aparece como
opção}\label{por-que-o-paraguai-aparece-como-opuxe7uxe3o}

Paraguai combina três características que explicam por que ele aparece
com frequência em discussões sobre autossuficiência e preparação: base
energética barata e abundante, baixa escala demográfica no interior do
país e uma estrutura econômica que ainda preserva espaço para produção
agropecuária e logística de menor atrito.

Em 2025, a ITAIPU informou o suprimento de 25.768 GWh ao Paraguai via
ANDE, com produção total de 72.879 GWh no ano, número equivalente a
cerca de 2 anos e 9 meses da demanda nacional. A Yacyretá também segue
como fonte estrutural de energia renovável e fornece o lastro de um
sistema elétrico altamente concentrado em hidrelétricas binacionais.
Para um público que busca resiliência, isso importa porque reduz o custo
estrutural de eletrificação, refrigeração, bombeamento e produção local.

O país também fechou 2025 com exportações totais de USD 16.720,3
milhões, puxadas principalmente por carne, milho e óleo de soja, e com o
regime de maquila em expansão. Isso mostra que a economia continua
ganhando tração, ainda que carregue a limitação de ser um país sem
litoral marítimo e, portanto, estruturalmente dependente de corredores
logísticos regionais.

Os pontos fracos existem e precisam ser explicitados. A maior parte da
população e dos serviços se concentra em Asunción e Central, o que
amplia a centralização de decisão, emprego e atendimento. Fora desse
eixo, o país oferece mais espaço e menor densidade, mas também exige
mais autossuficiência logística. Em outras palavras: Paraguai não é um
refúgio automático; é um país que pode funcionar bem para quem entende
suas assimetrias e escolhe o lugar certo.

\subsection{1.2 História, resiliência e
crescimento}\label{histuxf3ria-resiliuxeancia-e-crescimento}

O Paraguai é uma sociedade marcada por trauma histórico, reconstrução
lenta e adaptação contínua. Em registro oficial recente do próprio
governo, a trajetória nacional é descrita como uma história que
atravessou guerra, reconstrução, conflitos políticos e períodos de
reorganização institucional. Essa herança ajuda a explicar a cultura de
contenção, economia doméstica e resistência social que ainda aparece em
muitas comunidades.

Na leitura contemporânea do país, dois sinais são importantes. O
primeiro é a melhora gradual de indicadores sociais: a Presidência
informou que a pobreza monetária caiu para 20,1\% em 2024, o menor nível
desde a adoção da metodologia atual. O segundo é o avanço econômico
recente, com o BCP projetando crescimento do PIB de 6\% em 2025 e 4,2\%
em 2026, enquanto a inflação fechou 2025 em 3,1\%, já perto da meta
oficial.

Para o público-alvo do livro, a leitura prática é simples: trata-se de
um país que ainda cresce, ainda é relativamente acessível no interior e
ainda preserva espaço físico e institucional para projetos de médio
prazo. Isso não elimina os riscos, mas ajuda a entender por que o
Paraguai é percebido como ambiente de construção gradual e não apenas
como destino de crise.

\subsection{1.3 Recursos naturais e base
física}\label{recursos-naturais-e-base-fuxedsica}

O principal ativo hídrico subterrâneo que merece destaque é o Sistema
Aquífero Guarani. O MADES informou em 2025 que o recurso cobre 87.536
km² no território paraguaio e que o país vem ampliando a rede de
monitoramento com novos poços e estações automáticas. Isso importa
porque água subterrânea confiável é um dos pilares mais subestimados em
qualquer estratégia de permanência.

Na superfície, o Paraguai depende fortemente de grandes rios e da
hidrovía. O Rio Paraguai e o Rio Paraná estruturam navegação, fronteira,
produção e logística. O próprio MADES mantém operações de controle e
proteção nos rios Paraguay e Tebicuary, mostrando o peso desses
corredores hídricos na vida nacional. Em termos práticos, rios são
infraestrutura no Paraguai, não apenas paisagem.

Somando energia hidrelétrica, água subterrânea monitorada e vocação
agroexportadora, o país tem uma base física que pode favorecer
autossuficiência parcial ou progressiva, desde que a escolha territorial
seja feita com critério.

\subsection{1.4 Energia solar: o marco federal e o que varia por
concessionária}\label{energia-solar-o-marco-federal-e-o-que-varia-por-concessionuxe1ria}

O Paraguai construiu, entre 2023 e 2024, seu marco regulatório mais
abrangente para energias renováveis não convencionais. A \textbf{Lei N°
6977/\allowbreak 2023} regula o fomento, geração, produção, desenvolvimento e
utilização de energia elétrica a partir de fontes de energias renováveis
não convencionais não hidráulicas --- cobrindo solar fotovoltaica,
eólica, biogás, biomassa, geotérmica e hidrogênio verde. O
\textbf{Decreto N° 1168/\allowbreak 2024}, promulgado em fevereiro de 2024, detalha
os tipos de agentes reconhecidos, os processos de licenciamento e
consolida os incentivos fiscais: isenções de IVA e direitos
alfandegários para painéis fotovoltaicos, acumuladores de energia,
inversores e aerogeradores destinados a projetos registrados.

Para quem instala um sistema solar off-grid ou híbrido em propriedade
rural, o enquadramento relevante é como \textbf{Autogenerador ERNC}
perante o Viceministerio de Minas y Energía (MOPC). Abaixo de 1 MW de
capacidade instalada, o autogenerador opera com maior autonomia e sem
obrigação de fornecimento à rede. Para sistemas totalmente off-grid, o
registro é recomendável para acessar os benefícios fiscais; a ausência
de registro não gera sanção explícita para pequenos sistemas, mas impede
o acesso às isenções.

O que varia por concessionária ou departamento é relevante: a rede
nacional da ANDE cobre a Região Oriental e parte do Chaco, mas em zonas
não atendidas pela ANDE existem concessionárias ou cooperativas locais
com seus próprios procedimentos de aprovação e conexão para sistemas
grid-tied. Não há regulamento técnico unificado público entre todas as
concessionárias.

Em fevereiro de 2025, o Viceministerio anunciou um programa experimental
de autoconsumo solar residencial com previsão de licença específica ---
mas até março de 2026 os requisitos ainda estavam em elaboração,
configurando uma lacuna regulatória transitória para o segmento
residencial de pequeno porte.

\subsection{1.5 Infraestrutura de segurança: do federal ao
departamental}\label{infraestrutura-de-seguranuxe7a-do-federal-ao-departamental}

A arquitetura de segurança pública do Paraguai é coordenada pelo
\textbf{Ministério do Interior (MDI)}, que exerce a direção política da
Polícia Nacional e coordena os organismos de segurança. Abaixo do MDI
funcionam a \textbf{Polícia Nacional}, a \textbf{SENAD} (Secretaria
Nacional Antidrogas) e o sistema de emergências \textbf{911}. O
\textbf{MDN} (Ministério da Defesa) coordena as Forças Armadas, que
desde 2022-2023 passaram a ter papel ativo em segurança interna ---
especialmente nas fronteiras, com Forças de Tarefa Conjuntas (FTC Norte
no eixo Concepción-San Pedro-Amambay, e FTC Sul na fronteira com
Argentina).

A Polícia Nacional foi reorganizada pela \textbf{Lei N° 7280/\allowbreak 2024}
(Reforma e Modernização) e opera por meio de 18 jefaturas de polícia
departamentais. A distribuição real de efetivos é altamente desigual:
Central e Assunção concentram a maior parte; departamentos rurais como
Alto Paraguay, Boquerón e partes de Concepción têm cobertura muito mais
esparsa.

O \textbf{Sistema 911} (Lei N° 4739) opera com um Centro Nacional em
Assunção e centros regionais em Concepción, Encarnación, Coronel Oviedo,
Ciudad del Este e San Juan Bautista. Uma inovação relevante é o
\textbf{AML (Advanced Mobile Location)}, que envia automaticamente as
coordenadas GPS do chamador --- tornando o Paraguai o primeiro país da
América do Sul a obter licença para esta tecnologia. A rede de
videovigilância possui aproximadamente 1.641 câmeras instaladas, das
quais apenas \textasciitilde49\% estavam operacionais --- uma
fragilidade estrutural persistente.

A implicação prática para quem planeja se estabelecer no interior é
direta: em áreas rurais remotas e no Chaco, a presença real de segurança
pública pode se limitar a um posto policial isolado ou à presença
eventual das FFAA em patrulhamento de fronteira.

\subsection{1.6 Terra rural: tamanho mínimo, restrições e faixas de
preço}\label{terra-rural-tamanho-muxednimo-restriuxe7uxf5es-e-faixas-de-preuxe7o}

O marco legal fundamental é o \textbf{Estatuto Agrário} (Lei N°
1863/\allowbreak 2002), aplicado pelo \textbf{INDERT} (Instituto Nacional de
Desenvolvimento Rural e da Terra). O conceito central é a \textbf{UBEF
(Unidad Básica de Economía Familiar)}: a propriedade agrária necessária
para que uma família rural obtenha renda suficiente para sua fixação
efetiva. O Estatuto proíbe parcelamento em áreas menores que uma UBEF em
colônias agrícolas.

Tamanho mínimo na prática: o Estatuto fixou que, nos primeiros três anos
de vigência, a UBEF não seria inferior a \textbf{10 hectares}. Após esse
período transitório, o tamanho seria determinado por estudos técnicos do
INDERT por zona --- mas essa definição pós-transitória não foi
localizada em documento oficial publicado. A referência de 10 ha
permanece sendo o parâmetro prático mais usado no mercado por ausência
de atualização formal.

Acesso de estrangeiros: em geral, estrangeiros têm os mesmos direitos de
propriedade que cidadãos paraguaios. A exceção é a \textbf{zona de
segurança de fronteira} (Lei N° 2532/\allowbreak 2005) --- faixa de \textbf{50 km a
partir de todas as fronteiras terrestres e fluviais} com Argentina,
Brasil e Bolívia. Nessa faixa, estrangeiros dos países vizinhos
(argentinos, brasileiros e bolivianos) não podem ser proprietários de
imóveis rurais sem autorização do Poder Executivo. Imóveis urbanos
dentro da zona de fronteira não estão sujeitos a essa restrição.

Faixas de valor de terra rural por região (referência de mercado,
2024-2025):

{\def\LTcaptype{none} % do not increment counter
\begin{longtable}[]{@{}ll@{}}
\toprule\noalign{}
Região & Faixa de Valor (USD/\allowbreak ha) \\
\midrule\noalign{}
\endhead
\bottomrule\noalign{}
\endlastfoot
Alto Paraná (fronteira Brasil, premium) & 12.000 -- 15.000 \\
Caaguazú & 8.857 \\
Itapúa (fronteira Argentina) & 6.945 \\
San Pedro (norte da Região Oriental) & 4.608 \\
Concepción (norte) & 2.000 -- 2.500 \\
Região Oriental sul/\allowbreak centro (cerealeiro) & 3.400 -- 10.000 \\
Chaco central & 450 (virgem) -- 900 (formado) \\
Chaco geral & 500 -- 1.500 \\
\end{longtable}
}

Importante: esses valores são de mercado (plataformas privadas). O Valor
INDERT e o valor cadastral/\allowbreak fiscal são sistematicamente mais baixos ---
as três referências não coincidem, característica estrutural documentada
do mercado fundiário paraguaio.
